% =====================================================================
%  НЭГ. ТӨСЛИЙН ТАНИЛЦУУЛГА
% =====================================================================
\chapter{ТӨСЛИЙН ТАНИЛЦУУЛГА}

\section{Төслийн суурь үзүүлэлтүүд}

Төслийн суурь үзүүлэлтүүдийг Хүснэгт~\ref{tab:ch1-suur}~—т үзүүлэв.

\begin{table}[h]
  \centering
  \caption{Төслийн суурь үзүүлэлтүүд}
  \label{tab:ch1-suur}
  \begin{tabular}{@{}p{6.2cm}p{7.4cm}@{}}
    \toprule
    \textbf{Үзүүлэлт} & \textbf{Утга} \\ \midrule
    Төслийн нэр & «Ухаалаг кофе шопын NFC-д суурилсан захиалгын систем» \\
    Төсөл хэрэгжүүлэгч & М.Мөнхдорж, Ү.Тамир, Б.Эрдэнэбаяр \\
    Төслийн төрөл & Програм хангамжийн төсөл (бүтээгдэхүүний шууд
                      борлуулалттай) \\
    Хэрэгжих хугацаа & 6 сар (2026.10 -- 2027.03) \\
    Хөрөнгө оруулалтын хэмжээ & 55{,}533{,}896 төгрөг \\
    Төлөвлөж буй хүчин чадал & Нэг жилийн дотор 10 -- 20 кофе шопт
                               бүтээгдэхүүнээ борлуулах \\
    Цаашдын зорилго & Хот хөдөөгийн ялгаагүйгээр, улмаар гадаад зах зээлд
                      экспортлох \\
    \bottomrule
  \end{tabular}
\end{table}

\section{Төсөл хэрэгжүүлэгчийн тухай}

Төслийг Шинэ Монгол Технологийн Их Сургуулийн (ШМТИС) оюутнууд
М.Мөнхдорж, Ү.Тамир, Б.Эрдэнэбаяраас бүрдсэн баг бүтээн байгуулна. Баг
нь веб болон гар утасны аппликешн хөгжүүлэлт, өгөгдлийн сан, төслийн
удирдлагын чиглэлээр бэлтгэгдсэн гурван гишүүдээс бүрдэнэ. Багийн
гишүүд урьд өмнө веб болон гар утасны аппликешн хөгжүүлэх төслүүдэд
ажиллаж байсан туршлагатай, React, NestJS, PostgreSQL технологуудын
судалгааг тусгайлан хийсэн (\cite{react, nestjs, postgres}).

Санхүүгийн чадавх: Төслийн нийт хөрөнгө оруулалтын 50\% нь багийн
өөрийн санхүүгийн хөрөнгөөр, үлдсэн 50\% нь зээлийн хөрөнгөөр
хангагдана. Төслийн санхүүгийн урьдчилсан тооцоо нь хөрөнгө
оруулалтыг 35 сарын дотор нөхөхөд чиглэсэн бөгөөд нийлээд
санхүүгийн чадавх нь төслийг бүрэн хэрэгжүүлэхэд хүрэлцэхүйц болно.

\section{Төслийн тухай}

\subsection{Төслийн бүтээгдэхүүний танилцуулга}

Төслийн бүтээгдэхүүн нь кофе шопын ширээнд суусан хэрэглэгчид зориулсан
NFC стикер, PWA (Progressive Web App) маягийн цахим цэс, цахим
захиалга болон төлбөрийн системээс бүрдсэн цогц шийдэл юм. Хэрэглэгч
ширеэнд байрлах NFC стикерийг утсаар уншуулснаар тухайн ширээг
автоматаар тодорхойлж, цахим цэс нээгдэж, захиалгаа өгч, төлбөрийг
нутгаасаа төлснөөр захиалгаа ширээндээ хүргэж авна.

Төсөл хэрэгжүүлэгч баг үйлчилгээ үзүүлэгч бус, бүтээгдэхүүнийг шууд
борлуулах загвараар ажиллана. Ажилладаг кофе шоп бүрд нэг удаа
NFC стикерийг ширээ тус бүрт байршуулж, цэс болон холбогдох
системийг тухайн кофе шопт тохируулан өгснөөр манай ажил дуусна.
Дараа нь кофе шоп өөрийн админ эрхээр нэвтрэн цэсээ шинэчлэх --
шинэ бүтээгдэхүүн нэмэх, хасах, үнийн өөрчлөлт оруулах боломжтой.
Энэхүү загварыг олон кофе шоптой давтан хэрэгжүүлж, ашиг олно.

Бүтээгдэхүүний гол бүрэлдэхүүн хэсгүүд:

\begin{itemize}
  \item \textbf{NFC стикер (NTAG213/215):} Ширээ бүр дээр байрлах, тухайн
        ширээний дугаарыг кодчилсон стикер \cite{nxpntag}.
  \item \textbf{Хэрэглэгчийн PWA аппликешн:} Утасны арын дэвсгэртэд NFC
        уншиж, цэсийг нээх, захиалга өгөх, төлбөр төлөх интерфэйс \cite{androidnfc, corenfc}.
  \item \textbf{Кофе шопын админ хэсэг:} Цэс, үнийн сан, тайлан тооцоо,
        нэгжийн тохиргоог өөрөө удирдах бүрэн эрхт хэсэг.
  \item \textbf{Ажилтны хэсэг:} Захиалга хүлээн авах, гүйцэтгэх
        ажлын урсгалыг удирдах dashboard.
\end{itemize}

\subsection{Төслийн танилцуулга}

Уламжлалт кофе шопын үйлчилгээний урсгалд дараах бэрхшээлүүд
тархаж байна:

\begin{itemize}
  \item Кассын дээр хүн бөөгнөрч, үйлчлүүлэгчид удаан хүлээх.
  \item Төлбөрийн хэлбэр нь нэг удаагийн бөгөөд бичил болон бэлэн
        мөнгөний алдаа гарах боломжтой.
  \item Цэсээ хэвлэмэл хэлбэрээр харуулах нь өдөр тутмын хуучирал,
        засах зардлыг нэмэгдүүлдэг.
  \item Үйлчлүүлэгчийн мэдээллийг цуглуулж, үйлчилгээг хувь хүний
        тохирох боломжгүй.
\end{itemize}

Манай төсөл дээрх бэрхшээлүүдийг NFC технологи болон цахим төлбөрийн
системийг хослуулан шийдвэрлэнэ. Хэрэглэгчийн хөдөлгөөн дуусах хүртэлх
процессыг хурдасгаж, кофе шопын ажиллах хүчний ажлын ачааллыг бууруулж,
цаашид үйлчилгээний чанарыг нэмэгдүүлнэ.

\section{Төслийн үндсэн зорилго, зорилтууд}

\subsection{Төслийн зорилго, зорилтууд}

\textbf{Ерөнхий зорилго:} Кофе шопын үйлчилгээний урсгалыг NFC
технологиор автоматжуулан, хэрэглэгчид кассын дараалалгүйгээр захиалга
өгөх, төлөх, үйлчилгээг хүлээн авах нэгдсэн орчин бүрдүүлэх.

\textbf{Зорилтууд:}

\begin{enumerate}
  \item NFC стикерээр ширээг таних, цэсийг нэг алхамд нээх хэрэглэгчийн
        системийг бүрэн хэрэгжүүлэх.
  \item Захиалга өгөх, төлбөр төлөх, захиалгын төлөвийг хянах урсгалыг
        тасалдалгүйгээр холбох.
  \item Салбарын ажилтанд захиалгыг бодит цагт хүргэдэг хяналтын
        системийг боловсруулах.
  \item Бүрэн аюулгүй, хувийн мэдээллийг хамгаалсан архитектуртай
        системийг нэвтрүүлэх.
  \item Нэг жилийн дотор 10 -- 20 кофе шопт бүтээгдэхүүнээ
        борлуулж, системийн тогтвортой ажиллагааг батлах.
\end{enumerate}

\subsection{Төслийн ач холбогдол}

\begin{itemize}
  \item \textbf{Үйлчлүүлэгчид:} Хүлээх хугацааг 40 -- 50\% бууруулж,
        кассын дараалалгүй, нэг аппаар захиалга болон төлбөрийг
        гүйцэтгэх боломж.
  \item \textbf{Байгууллагад:} Кассын алдаа буурч, хүний нөөцийн
        ачаалал 20 -- 30\% хэмнэгдэнэ. Үйлчлүүлэгчийн мэдээлэл
        цуглуулагдаж, маркетингийн шийдвэр гаргалт илүү үр дүнтэй
        болно.
  \item \textbf{Салбарт:} Дотоодын компьютерийн технологийн
        бүтээгдэхүүний хомсдолыг нөхөж, уламжлалт бизнесийн салбарт
        цифржилтийг түлхүүлэх нөлөөтэй.
  \item \textbf{Эдийн засагт:} Уламжлалт үйлчилгээний салбарын
        урт хугацааны хэрэгцээнд чиглэсэн ажлын байр, технологийн
        шийдлийг дотоодын хүчээр бүтээх чадварыг нэмэгдүүлнэ.
\end{itemize}
